% =============================================================================
% CHAPTER EXTENSIONS (SET II): Ge-on-Si PIN Photodetector
% Islam I. Abdulaal, 2026
%
% SOURCE REVIEWS ANALYSED:
%   [A] X.-X. Wang et al., "Photodetectors integrating waveguides and
%       semiconductor materials," Nanoscale, vol. 16, pp. 5504–5520, 2024.
%       [cited as \cite{Wang2024}]
%   [B] T. Yan et al., "High-performance Ge photodetectors on silicon
%       photonics platform for optical interconnect," Sensors and Actuators A,
%       vol. 376, p. 115535, 2024.
%       [cited as \cite{Yan2024}]
%
% INTEGRATION INSTRUCTIONS:
%   Each block is labelled with its target section and a precise anchor
%   sentence (the last line of the paragraph it follows). Paste each
%   extension immediately after that anchor, before the next structural
%   heading. Do NOT alter existing text. Both \cite{Wang2024} and
%   \cite{Yan2024} are already present in the provided .bib file.
% =============================================================================


% -----------------------------------------------------------------------------
% EXTENSION R1
% Target section : \subsubsection{Responsivity and Quantum Efficiency}
%                  (\cref{subsubsec:pd_responsivity})
% Insert after   : "...confirming that the absorber length $L$ is the primary
%                  responsivity design lever through the Beer-Lambert factor of
%                  \cref{eq:qe_chain}."
% -----------------------------------------------------------------------------

A more complete decomposition of the quantum efficiency chain inserts an
additional factor that is suppressed in \cref{eq:qe_chain} whenever the device
is fully depleted and all photogenerated carriers drift at saturation velocity,
but that becomes quantitatively significant under partial depletion or at zero
bias. Denoting by $\xi\in(0,1]$ the fraction of absorbed photons whose
electron-hole pairs are actually collected at the external contacts—as opposed
to recombining within the absorber before reaching the depletion boundary—the
external quantum efficiency takes the more general form~\cite{Yan2024}:
\begin{equation}
    \eta_{\mathrm{ext}}
    = (1 - R_{\mathrm{refl}})\,\xi\,\eta_c\,
      \bigl(1 - e^{-\alpha(\lambda)\,L}\bigr),
    \label{eq:qe_full}
\end{equation}
where $R_{\mathrm{refl}}$ is the power reflectance at the Si/Ge entry facet of
the taper. The parameter $\xi$ is not a material constant but a function of
the local electric field: in the fully depleted intrinsic region at saturation
field, the drift transit time $\tau_{\mathrm{tr}} = W_i/v_{\mathrm{d,sat},h}$
is much shorter than the SRH recombination lifetime $\tau_0$, giving
$\xi \approx 1 - \tau_{\mathrm{tr}}/\tau_0 \approx 1$ to an excellent
approximation for the parameters of \cref{tab:layer_stack}. It is precisely
this condition—$\tau_{\mathrm{tr}} \ll \tau_0$—that makes the IQE of
\cref{eq:iqe} approach unity and that justifies the simpler form of
\cref{eq:qe_chain} as the operational limit of \cref{eq:qe_full}. Any bias
reduction toward zero, however, lowers the electric field, slows the drift
velocity below its saturated value, and allows $\xi$ to depart measurably from
unity, thereby degrading the photocurrent even when the total absorbed power
$P_{\mathrm{abs}}$ remains unchanged. The factored form of \cref{eq:qe_full}
is therefore the natural framework within which zero-bias or low-bias operating
points—motivated by the energy-per-bit metric $e_c \propto |V_{\mathrm{bias}}|
\,\mathcal{R}$ introduced in \cref{sec:pd_discussion}—must be evaluated: both
$\xi$ and $\eta_c$ may degrade simultaneously as the bias is relaxed, and their
combined reduction cannot be identified from the responsivity alone without an
independent measurement of the absorbed power $P_{\mathrm{abs}}$.


% -----------------------------------------------------------------------------
% EXTENSION R2
% Target section : \subsection{Band Structure of Germanium and the Direct-Gap
%                  Absorption Mechanism} (\cref{subsec:pd_bandstructure})
% Insert after   : "...cannot be used as the absorber. This single
%                  material-physics argument is the entire justification for
%                  the Ge-on-Si heterostructure that forms the subject of this
%                  chapter."
%   [Insert after the same anchor as Extension 2, but place AFTER Extension 2]
% -----------------------------------------------------------------------------

A further consequence of the Ge band structure, not relevant to the nominal
operating point of this device but important to the design space of next-generation
detectors, is the sensitivity of the direct-gap energy to biaxial in-plane
tensile strain $\varepsilon$. The application of biaxial tensile strain lifts
the heavy-hole band above the light-hole band at $\Gamma$ and reduces the
$\Gamma$-valley conduction band minimum relative to the $L$-valley, modifying
the direct gap according to the linear deformation-potential relation:
\begin{equation}
    E_{g,\mathrm{dir}}(\varepsilon)
    \approx E_{g,\mathrm{dir}}(0)
      - \bigl(2a_v + a_c\bigr)\bigl(1 - c_{12}/c_{11}\bigr)\varepsilon,
    \label{eq:bandgap_strain}
\end{equation}
where $a_v$ and $a_c$ are the valence- and conduction-band hydrostatic
deformation potentials, and $c_{11}$, $c_{12}$ are the elastic stiffness
constants of germanium~\cite{Yan2024}. The net effect for the biaxial tensile
strain of approximately $0.2$--$0.3\%$ that develops in Ge epitaxial layers
during the post-growth cooldown from the high-temperature deposition step—owing
to the thermal expansion coefficient mismatch between Ge and Si—is a reduction
of $E_{g,\mathrm{dir}}$ by approximately \SI{25}{\milli\electronvolt} relative
to the unstrained bulk value of \SI{0.80}{\electronvolt}. This shift red-shifts
the onset of direct-gap absorption by roughly \SI{50}{\nano\metre}, extending
the high-absorption edge of the intrinsic layer slightly beyond the nominal
\SI{1550}{\nano\metre} C-band boundary and enhancing $\kappa(\lambda)$ at
C-band wavelengths relative to the bulk tabulation used in
\cref{eq:alpha_kappa}. For the O-band operating point of the present device
the strain correction is small and absorbed within the uncertainty of the
tabulated optical constants; it is, however, the principal physical mechanism
exploited in strained-Ge emitter and extended-wavelength detector designs, and
its formal origin in \cref{eq:bandgap_strain} makes clear that it represents a
controlled modification of the same direct-gap absorption edge—at the
$\Gamma$-point—that governs the primary absorption process described in
\cref{subsec:pd_bandstructure}.


% -----------------------------------------------------------------------------
% EXTENSION R3
% Target section : \subsubsection{Dark Current}
%                  (\cref{subsubsec:pd_dark})
% Insert after   : "...and explains why SRH generation dominates in reverse-
%                  biased Ge junctions at room temperature, where
%                  $n_i^2/N_{D,A} \ll n_i$."
% -----------------------------------------------------------------------------

A physically distinct and practically important contribution to the total dark
current, not captured by the bulk generation terms of \cref{eq:idark_total},
arises from the finite density of interface trap states at the Ge surface in
contact with the passivation oxide. These traps, located within the forbidden
gap at the Ge/SiO$_2$ interface, act as generation centres governed by an
interface-specific SRH rate whose magnitude is proportional to the interface
trap density $D_{it}$ (in units of \si{\per\electronvolt\per\centi\metre\squared})
and the interface recombination velocity defined in \cref{eq:jdark_surf}. In
practice, the dark current density of a Ge photodiode partitions into bulk and
surface contributions as:
\begin{equation}
    I_d = I_{\mathrm{bulk}} + I_{\mathrm{surf}}
    = A_j \bigl(J_{\mathrm{SRH}} + J_{\mathrm{diff}}\bigr)
    + P_j\,q\,n_i\,v_s/2,
    \label{eq:idark_decomposed}
\end{equation}
where $P_j = 2(W_m + L)$ is the perimeter of the junction mesa and the second
term represents the surface leakage current driven by the perimeter-normalised
recombination velocity $v_s$~\cite{Yan2024}. This perimeter scaling has a
direct and non-trivial geometric consequence: as the junction area $A_j$ is
scaled down to reduce capacitance $C_j$ in \cref{eq:cj}—which is a
monotonically beneficial modification for bandwidth—the area-to-perimeter ratio
$A_j/P_j = W_m L/[2(W_m + L)]$ decreases, so the fractional contribution of
the surface term in \cref{eq:idark_decomposed} increases. For the present
device with $A_j = \SI{6}{\micro\metre\squared}$ and $P_j = \SI{14}{\micro\metre}$,
the area-to-perimeter ratio is $A_j/P_j \approx \SI{0.43}{\micro\metre}$; a
future device scaled to $A_j = \SI{2}{\micro\metre\squared}$ and correspondingly
narrower mesa dimensions would reduce this ratio toward \SI{0.17}{\micro\metre},
at which point surface-leakage suppression through Ge$_x$O$_y$ or Al$_2$O$_3$
passivation—which reduces $D_{it}$ and thereby $v_s$ by one to two orders of
magnitude—becomes a primary design lever independent of and complementary to
the SRH lifetime engineering represented by $\tau_0$. The minimum achievable
dark current density in Ge-on-Si epitaxial material is not set by device
geometry alone but also by the threading dislocation density (TDD) intrinsic
to the heteroepitaxial film: each threading dislocation segment provides a
spatially extended generation centre whose contribution to $J_{\mathrm{SRH}}$
scales approximately linearly with TDD, so that TDD $\approx \SI{5e7}{\per
\centi\metre\squared}$—the value achievable with the standard two-step
low-temperature/high-temperature growth sequence—corresponds to a bulk dark
current density floor of order \SI{10}{\milli\ampere\per\centi\metre\squared}
at \SI{-1}{\volt}~\cite{Yan2024}. The calibrated $\tau_0 = \SI{5}{\nano\second}$
in \cref{subsec:pd_charge} implicitly reflects this dislocation density, and
any order-of-magnitude improvement in TDD—achievable through aggressive in-situ
annealing or arsenic-assisted dislocation annihilation—would proportionally
reduce $J_{\mathrm{SRH}}$ and the corresponding shot noise floor of
\cref{eq:shot}, with direct implications for the sensitivity expression
\cref{eq:sens_full}.


% -----------------------------------------------------------------------------
% EXTENSION R4
% Target section : \subsubsection{Electrical Bandwidth}
%                  (\cref{subsubsec:pd_bw})
% Insert after   : "...the principal motivation for the U-shaped electrode
%                  architecture discussed in \cref{subsec:pd_architecture}."
%   [i.e., at the end of the existing bandwidth subsection]
% -----------------------------------------------------------------------------

It is instructive to derive explicitly the intrinsic-layer thickness at which
the transit-time and RC contributions to \cref{eq:bw_total} are equal in
magnitude, since this point defines the geometry at which further single-variable
optimisation of $W_i$ ceases to improve $f_{3\,\mathrm{dB}}$. Setting
$f_{\mathrm{tt}}(W_i^*) = f_{\mathrm{RC}}(W_i^*)$ and substituting
\cref{eq:bw_tt,eq:bw_rc,eq:cj}:
\begin{equation}
    \frac{0.45\,v_{\mathrm{d,sat},h}}{W_i^*}
    = \frac{W_i^*}{2\pi R_{\mathrm{tot}}\,\varepsilon_0\varepsilon_{\mathrm{Ge}}\,A_j},
    \label{eq:wi_optimal_eq}
\end{equation}
which yields the optimal thickness:
\begin{equation}
    W_i^* = \sqrt{
      \frac{0.45\,v_{\mathrm{d,sat},h}\cdot
            2\pi R_{\mathrm{tot}}\,\varepsilon_0\varepsilon_{\mathrm{Ge}}\,A_j}{1}
    }
    = \sqrt{0.9\pi\,R_{\mathrm{tot}}\,\varepsilon_0\varepsilon_{\mathrm{Ge}}
             \,v_{\mathrm{d,sat},h}\,A_j}.
    \label{eq:wi_optimal}
\end{equation}
Substituting $R_{\mathrm{tot}} = \SI{50}{\ohm}$, $\varepsilon_{\mathrm{Ge}} =
16.2$, $v_{\mathrm{d,sat},h} = \SI{4e6}{\centi\metre\per\second}$, and the
junction area $A_j = \SI{6}{\micro\metre\squared}$ of the present device gives
$W_i^* \approx \SI{310}{\nano\metre}$, within \SI{11}{\percent} of the chosen
$W_i = \SI{350}{\nano\metre}$—confirming that the device is positioned near
the bandwidth optimum consistent with the $R_{\mathrm{tot}} = \SI{50}{\ohm}$
standard load~\cite{Yan2024}. Any reduction of $R_{\mathrm{tot}}$ through
co-design with a lower-impedance transimpedance amplifier would shift $W_i^*$
downward, allowing a thinner absorber that simultaneously increases
$f_{\mathrm{tt}}$ and $f_{\mathrm{RC}}$ until the new balance point is
reached. The expression \cref{eq:wi_optimal} also exposes a fundamental scaling
law: since $W_i^* \propto \sqrt{A_j}$, shrinking the junction area by a factor
of four (reducing $W_m$ or $L$ by two) allows $W_i^*$ to be halved, raising
the balanced bandwidth at the new optimum by approximately the same factor. This
is the physical basis for the observation that the \SI{265}{\giga\hertz}
lateral-junction device of Lischke~\textit{et al.}~\cite{Lischke2021} requires
an absorber only \SI{90}{\nano\metre} thick: the dramatically reduced $A_j$ of
that architecture shifts $W_i^*$ far below the values practical in the vertical
architecture, where $A_j$ is constrained from below by the optical absorption
requirement of \cref{eq:abs_fraction}.

An additional caveat on the validity of \cref{eq:bw_total} deserves explicit
statement. The formula combines $f_{\mathrm{tt}}$ and $f_{\mathrm{RC}}$ as
independent roll-off poles, an approximation that is exact only when the carrier
impulse response is strictly rectangular (uniformly distributed generation and
purely drift transport). Carriers generated outside the depletion region—in the
heavily doped contact layers—contribute a diffusion tail to the impulse
response that is not rectangular and that introduces a low-frequency component
in the photocurrent spectrum, reducing the effective $f_{3\,\mathrm{dB}}$ below
the prediction of \cref{eq:bw_total}~\cite{Yan2024}. For the present device,
the $n^{++}$-Ge cathode is limited to \SI{50}{\nano\metre} specifically to
suppress this diffusion contribution, and the heavy doping of the contact layer
ensures that its minority-carrier diffusion length $L_{p,\mathrm{cat}}
= \sqrt{D_p\tau_p}$ is much shorter than its physical thickness, so that
essentially no photogeneration occurs within it at \SI{1310}{\nano\metre}.
This design choice is therefore not only motivated by free-carrier absorption
reduction through \cref{eq:qe_chain} but simultaneously necessary for the
validity of the transit-time bandwidth model of \cref{eq:bw_tt}.


% -----------------------------------------------------------------------------
% EXTENSION R5
% Target section : \subsection{Evolution of Waveguide-Integrated Ge-on-Si
%                  Photodetectors} (\cref{subsec:pd_history})
% Insert after   : "...without any modification of the junction geometry."
%   [i.e., after the Novack et al. inductive peaking paragraph]
% -----------------------------------------------------------------------------

It is worth situating the inductive gain-peaking strategy of
Novack~\textit{et al.}~\cite{Novack2013} within a broader equivalent-circuit
taxonomy, because the three canonical peaking topologies—series gain peaking,
enhanced series gain peaking, and shunt peaking—produce qualitatively different
modifications to the small-signal transfer function and carry different
implementation constraints in a photonic integrated circuit
context~\cite{Yan2024}. In series gain peaking, an inductance is inserted in
series with the load resistance $R_L$, forming an LC resonance with the junction
capacitance $C_j + C_{\mathrm{par}}$ that elevates the frequency response in
the vicinity of the resonance frequency $\omega_0 = 1/\sqrt{L(C_j +
C_{\mathrm{par}})}$; the achievable bandwidth extension relative to the passive
$f_{\mathrm{RC}}$ of \cref{eq:bw_rc} depends on the $Q$-factor of the inductor
and can theoretically reach a factor of two for ideal reactive elements, but
the distributed parasitic resistance of an on-chip spiral inductor limits the
practical improvement to approximately $40$--$60\%$. The shunt-peaking variant
places the inductance in parallel with $R_L$, shifting the pole structure of
the transfer function and providing a smoother, Butterworth-like extension of
the passband without the resonant overshoot of the series topology. In either
case, the gain-peaking element transforms the photodetector from a low-pass
RC device to a bandpass or flat-extended device, and the resulting mismatch
between the inductive source impedance and the capacitive input of a
conventional transimpedance amplifier (TIA) must be absorbed into the
co-design of the detector-TIA interface. This constraint becomes increasingly
stringent as the photodetector is scaled toward sub-femtofarad capacitance,
because the inductive impedance required to cancel $C_j$ at \SI{100}{\giga\hertz}
is only a few tens of picohenries—comparable to the parasitic inductance of
bond wires—making the distinction between intentional and parasitic inductance
ambiguous. The U-shaped electrode architecture of the present device deliberately
avoids this co-design complexity by achieving $f_{\mathrm{RC}} > f_{\mathrm{tt}}$
through passive $R_s$ reduction alone, maintaining a purely resistive source
impedance that is straightforwardly matched to a TIA.


% -----------------------------------------------------------------------------
% EXTENSION R6
% Target section : \subsection{Comparative Summary}
%                  (\cref{subsec:pd_compare})
% Insert after   : "...without modifications to the absorber material,
%                  thickness, or doping profile."
%   [i.e., at the end of the comparative summary section]
% -----------------------------------------------------------------------------

Two architectural alternatives to the vertical and lateral PIN topologies
documented in \cref{tab:pd_literature} provide additional context for the
design choices made in the present work and define the boundaries of the design
space more precisely. The first is the resonant-cavity-enhanced (RCE) vertical
photodiode, in which the absorber is placed within a Fabry-Pérot resonator
formed by a high-reflectivity distributed Bragg reflector (DBR) at the substrate
interface and a partial reflector at the top surface. At resonance, the
single-pass absorbed fraction $\mathcal{A}$ of \cref{eq:abs_fraction} is
replaced by the cavity-enhanced expression:
\begin{equation}
    \mathcal{A}_{\mathrm{RCE}}
    = \frac{(1 - R_1)(1 + R_2\,e^{-\alpha L})^2
            \left(1 - e^{-2\alpha L}\right)}
           {1 - 2\sqrt{R_1 R_2}\,e^{-\alpha L}\cos(4\pi n L/\lambda)
            + R_1 R_2\,e^{-2\alpha L}},
    \label{eq:rce_absorption}
\end{equation}
where $R_1$ and $R_2$ are the top and bottom mirror reflectances and $n$ is the
real refractive index of the absorber~\cite{Yan2024}. At the resonance condition
$4\pi n L/\lambda = 2m\pi$ ($m$ integer), \cref{eq:rce_absorption} reduces to
$\mathcal{A}_{\mathrm{RCE}} = (1-R_1)(1+\sqrt{R_2})^2/(1-\sqrt{R_1
R_2})^2$ for $e^{-\alpha L} \approx 1$, which can substantially exceed unity
times the single-pass value and enables near-unity quantum efficiency in absorbers
as thin as \SI{200}{\nano\metre}—a geometry that simultaneously yields
capacitances of order \SI{1}{\femto\farad} and transit-time bandwidths
approaching \SI{60}{\giga\hertz}~\cite{Yan2024}. The fundamental limitation of
the RCE architecture is its extreme wavelength selectivity: the resonance
bandwidth scales as $(1 - \sqrt{R_1 R_2})/\sqrt{R_1 R_2}$ multiplied by the
free spectral range, so that achieving $\mathcal{A}_{\mathrm{RCE}} \gg
\mathcal{A}$ requires high mirror reflectances that simultaneously narrow the
usable optical bandwidth to a few nanometres—incompatible with the
temperature-insensitive broadband operation required of data-centre receivers
where the junction temperature may fluctuate by tens of kelvins. The second
alternative is the uni-travelling-carrier (UTC) vertical structure, in which
only electrons contribute to the drift current: a heavily $p$-doped Ge
absorption layer is followed by a wider-bandgap (Si) depletion layer, so that
photogenerated holes relax quasi-instantaneously within the absorber while
electrons—with their higher saturation velocity $v_{\mathrm{d,sat},e} =
\SI{6e6}{\centi\metre\per\second}$ rather than the \SI{4e6}{\centi\metre\per
\second} of holes—traverse the depletion layer as the sole drift carriers. The
transit-time limit of \cref{eq:bw_tt} is accordingly replaced by
$f_{\mathrm{tt,UTC}} = 0.45\,v_{\mathrm{d,sat},e}/W_{\mathrm{dep}}$, which
for the same depletion width yields a transit-time bandwidth approximately
$50\%$ higher than the hole-limited value. The UTC design simultaneously
suppresses the space charge effect at high optical power by eliminating the
slower hole contribution to carrier density buildup within the depletion
region—the mechanism responsible for the bandwidth degradation under strong
illumination identified in \cref{sec:pd_discussion}—and thereby raises
$P_{\mathrm{screening}}$ of \cref{eq:p_screening} for a given reverse
bias~\cite{Yan2024}. Its principal fabrication challenge is the requirement
for a diffusion-blocking $p^+$-Ge/$n$-Si heterojunction that confines holes
to the absorber without creating a conduction-band barrier for electrons, a
constraint that has been met in silicon photonics only by careful doping profile
engineering at the Ge/Si heterointerface.


% -----------------------------------------------------------------------------
% EXTENSION R7
% Target section : \subsection{Device Architecture and Design Rationale}
%                  (\cref{subsec:pd_architecture})
% Insert after   : "...as verified by the FDTD simulation of
%                  \cref{subsec:pd_fdtd}."
% -----------------------------------------------------------------------------

The choice of evanescent wave coupling, as opposed to the direct butt-coupling
geometry in which the Si waveguide facet abuts the Ge absorber end-face, deserves
a more detailed physical justification in the context of the generation-rate
distribution $G_{\mathrm{opt}}(\mathbf{r})$ of \cref{eq:gen_rate}. In the
butt-coupled geometry, the full waveguide modal power is launched directly into
the Ge absorber end-face, producing a strongly peaked $G_{\mathrm{opt}}$ at the
longitudinal entry plane ($z = 0$) that decays exponentially over the absorption
length $\alpha^{-1} \approx \SI{4.8}{\micro\metre}$. This front-loading
concentrates photogenerated carriers at a single cross-section of the absorber,
inducing a local carrier density that can substantially screen the internal
electric field under moderate optical powers—precisely the mechanism described
by \cref{eq:p_screening}—and thereby degrading bandwidth at the input aperture
before the optical power has been significantly absorbed. Evanescent wave
coupling mitigates this concentration by distributing the energy transfer over
the full length of the taper, which in the present device spans \SI{40}{\micro
\metre}: the optical mode in the Si waveguide progressively transfers its power
to the Ge absorber as the modal overlap between the two waveguiding regions
increases along the taper length, producing a $G_{\mathrm{opt}}$ profile that
is more axially distributed than the single-pass Beer-Lambert exponential and
that maintains lower peak carrier density for a given total absorbed power.
Furthermore, evanescent coupling decouples the optical propagation direction
from the direction of carrier collection: photogenerated carriers travel
vertically from the intrinsic Ge layer to the ohmic contacts, orthogonal to
the longitudinal optical propagation, so that the device length $L$ can be
independently optimised for absorption efficiency through \cref{eq:abs_fraction}
without affecting the vertical drift time $W_i/v_{\mathrm{d,sat},h}$ that
enters \cref{eq:bw_tt}~\cite{Wang2024,Yan2024}. This orthogonality—the defining
feature of waveguide-coupled over surface-coupled photodetectors—is what enables
the simultaneous optimisation of $\mathcal{R}$ through $L$ and of
$f_{\mathrm{tt}}$ through $W_i$, a degree of freedom entirely absent in
surface-illuminated (normal-incidence) vertical photodetectors where the
thickness of the absorber simultaneously determines the absorption and the
transit time.


% -----------------------------------------------------------------------------
% EXTENSION R8
% Target section : \subsection{Optical Simulation: Three-Dimensional FDTD}
%                  (\cref{subsec:pd_fdtd})
% Insert after   : "The absorber length is swept from \SIrange{2}{8}{\micro
%                  metre} to identify the saturation point of $\mathcal{A}(L)$
%                  from \cref{eq:abs_fraction}."
% -----------------------------------------------------------------------------

The saturation of $\mathcal{A}(L)$ identified in the sweep is not a smooth
asymptote when the waveguide loss of the silicon input guide is taken into
account. For a Si ridge waveguide with propagation loss $\gamma_{\mathrm{wg}}
\approx \SI{3.6}{\decibel\per\centi\metre}$ at \SI{1310}{\nano\metre}, the
power delivered to the taper entry is attenuated by a factor
$\exp(-\gamma_{\mathrm{wg}} \ell_{\mathrm{wg}}/4.343)$ relative to the chip
input facet, where $\ell_{\mathrm{wg}}$ is the length of the routing waveguide~\cite{Wang2024}. This waveguide insertion loss is spectrally flat to first
order and does not affect the shape of $\mathcal{A}(L)$, but it reduces
$P_{\mathrm{opt}}$ in \cref{eq:responsivity} and therefore the responsivity
as measured at the chip input rather than at the taper entry. For a routing
length of $\ell_{\mathrm{wg}} = \SI{1}{\milli\metre}$, the correction is
approximately \SI{0.36}{\decibel}, equivalent to a responsivity reduction
of about \SI{4}{\percent} relative to the taper-input value of
$\mathcal{R} = \SI{0.95}{\ampere\per\watt}$. While this correction is small
for the present device, it becomes the dominant extrinsic loss mechanism in
photonic integrated circuits where the routing distance between the fibre-chip
coupler and the detector exceeds several millimetres, and it sets a practical
upper bound on the usable waveguide length for a given link power budget. The
FDTD simulation domain of the present work terminates at the taper entry, so
the simulated $\mathcal{R}$ corresponds to the taper-input responsivity; any
comparison with chip-level measurements must account for the routing waveguide
attenuation separately through the factor $10^{-\gamma_{\mathrm{wg}}
\ell_{\mathrm{wg}}/10}$ applied to \cref{eq:responsivity}.


% -----------------------------------------------------------------------------
% EXTENSION R9
% Target section : \section{Discussion} (\cref{sec:pd_discussion})
% Insert as a new \paragraph immediately after the existing paragraph titled
%   "Comparison with the literature and architectural constraints."
% -----------------------------------------------------------------------------

\paragraph{Nano-optical absorption enhancement as a path beyond the Beer-Lambert
limit.}
The present design operates within the single-pass Beer-Lambert framework of
\cref{eq:beer_lambert}, in which the absorption fraction $\mathcal{A}(L)$
is a monotonically saturating function of $L$ and can never reach unity for
a finite absorber at the bulk absorption coefficient $\alpha$. Several
nano-optical architectures transcend this single-pass ceiling by engineering
the electromagnetic boundary conditions within the absorber to increase the
effective interaction length beyond the physical thickness. Grating-guided-mode
resonance (GMR) structures introduce a sub-wavelength periodic perturbation
on the Ge surface that diffracts normally incident light into guided modes
propagating laterally within the absorber film; the coupling of these guided
modes to the vertical free-space continuum produces a Fano-type spectral
resonance in the absorption spectrum, and at the resonance condition the
effective optical path length within a \SI{200}{\nano\metre} Ge film can equal
that of a bulk film more than thirty times thicker~\cite{Yan2024}. The physical
condition for maximum GMR absorption in a film of thickness $d$ and refractive
index $n$ is that the grating period $\Lambda$ satisfies
$\Lambda \approx \lambda/(n_{\mathrm{eff},m} - \sin\theta_i)$, where
$n_{\mathrm{eff},m}$ is the effective index of the $m$-th guided mode and
$\theta_i$ is the incidence angle. For the present device geometry the
waveguide-coupled field is not normally incident but propagates longitudinally;
the GMR concept is therefore more directly applicable to the normal-incidence
(surface-illuminated) architecture, where the trade-off between absorber
thickness and transit time is most acute, than to the waveguide-coupled
architecture studied here. Nevertheless, the GMR mechanism identifies a
physical principle—increasing the effective interaction length within a fixed
absorber thickness through resonant mode coupling—that is structurally analogous
to the photonic crystal (PhC) hole arrays that have been integrated into
waveguide Ge PDs to achieve broadband absorption enhancement over the full
C-band from a single \SI{350}{\nano\metre} absorber at the cost of increased
dark current from the etched hole sidewalls~\cite{Yan2024}. For the vertical
n-i-p architecture of the present chapter, the geometric constraint that
$A_j = W_m \times L$ must remain large enough for the Beer-Lambert absorption
to saturate in a \SI{6}{\micro\metre} length acts as a lower bound on both
$C_j$ through \cref{eq:cj} and on the dark current through
\cref{eq:jdark_srh}; a GMR or PhC overlay that achieves the same $\mathcal{A}$
in a shorter $L$—or equivalently in a thinner $W_i$—would therefore allow
$A_j$ to be reduced, simultaneously decreasing $C_j$ and the dark current
generation volume without sacrificing responsivity, and would push $W_i^*$ of
\cref{eq:wi_optimal} toward smaller values consistent with higher balanced
bandwidths. Whether the fabrication complexity and sidewall-defect-induced dark
current increase of etched nano-optical structures outweigh these benefits in a
specific process node is a quantitative question that cannot be resolved by the
present simulation framework but that represents a concrete and physically
well-posed direction for future device optimisation.

\paragraph{High-power operating limit and the space-charge saturation current.}
The analysis of \cref{sec:pd_results} assumes linear operation, in which the
photocurrent is proportional to the incident optical power and the electric
field distribution within the intrinsic layer is unperturbed by the
photogenerated carrier density. This linearity assumption fails when the
free-carrier charge density $\rho_{\mathrm{ph}} = q(n_{\mathrm{ph}} +
p_{\mathrm{ph}})$ generated by the optical field becomes comparable in magnitude
to the background ionised doping density, partially screening the built-in and
applied electric fields. The threshold optical power above which this
space-charge screening degrades the small-signal bandwidth by more than
\SI{1}{\decibel}—the so-called $-1$\,dB compression power $P_{\mathrm{-1dB}}$—
can be estimated by requiring $\rho_{\mathrm{ph}} < q N_D^+$ within the
depletion region, which translates to a constraint on the photocurrent
density~\cite{Yan2024}:
\begin{equation}
    J_{\mathrm{ph}} = \mathcal{R}\,P_{\mathrm{opt}}/A_j
    < q\,v_{\mathrm{d,sat},h}\,N_{\mathrm{eff}},
    \label{eq:saturation_current_density}
\end{equation}
where $N_{\mathrm{eff}}$ is an effective doping level that depends on the
junction profile. For the present device, substituting
$v_{\mathrm{d,sat},h} = \SI{4e6}{\centi\metre\per\second}$ and assuming
$N_{\mathrm{eff}} \sim \SI{e13}{\per\centi\metre\cubed}$ (the intrinsic carrier
density, as the absorber is nominally undoped) gives a limiting photocurrent
density of approximately \SI{64}{\milli\ampere\per\centi\metre\squared},
corresponding to a saturation photocurrent of $I_{\mathrm{sat}} \approx
J_{\mathrm{ph,sat}} \times A_j \approx \SI{3.8}{\micro\ampere}$ for the
\SI{6}{\micro\metre\squared} junction—in order-of-magnitude agreement with the
onset of field screening at $P_{\mathrm{opt}} \approx P_{\mathrm{screening}}$
estimated in \cref{sec:pd_discussion}. Operating at the \SI{100}{\micro\watt}
input power of the present simulation gives a photocurrent of
$I_{\mathrm{ph}} = \SI{95}{\micro\ampere}$ at $\mathcal{R} =
\SI{0.95}{\ampere\per\watt}$, which corresponds to a photocurrent density of
\SI{1.58}{\milli\ampere\per\centi\metre\squared}—well below the saturation
threshold and confirming that the linear assumption is justified at the nominal
operating point. In contrast, for data-centre receiver applications where the
received optical power can reach $-3\,\mathrm{dBm} \approx \SI{500}{\micro\watt}$
after waveguide routing losses, the photocurrent density would approach
\SI{7.9}{\milli\ampere\per\centi\metre\squared}, which begins to enter the
regime where space-charge screening must be accounted for in the bandwidth
prediction. The UTC architecture discussed in \cref{subsec:pd_compare}
circumvents this limitation by using only the faster electron species as the
drift carrier, raising the right-hand side of \cref{eq:saturation_current_density}
by the ratio $v_{\mathrm{d,sat},e}/v_{\mathrm{d,sat},h} = 1.5$, and by
operating with a heavily $p$-doped absorber that increases $N_{\mathrm{eff}}$
above the intrinsic value, thereby extending the linear dynamic range and the
compression power before any reduction in $f_{3\,\mathrm{dB}}$ is observed.


% -----------------------------------------------------------------------------
% EXTENSION R10
% Target section : \section{Device-to-System Projection} (\cref{sec:pd_system})
% \subsection{Bandwidth and Modulation Format Implications}
%                  (\cref{subsec:pd_format})
% Insert after   : "...provides the input channel model for that equaliser
%                  design."
% -----------------------------------------------------------------------------

The modulation-format implications of the photodetector channel model can be
quantified more precisely by evaluating the inter-symbol interference (ISI)
penalty imposed by the Ge PD transfer function on the received signal. Modelling
the photodetector as a first-order low-pass filter with $-3\,\mathrm{dB}$
frequency $f_{3\,\mathrm{dB}} = \SI{103}{\giga\hertz}$ and flat-band
responsivity $\mathcal{R} = \SI{0.95}{\ampere\per\watt}$, the attenuation of a
sinusoidal component at the Nyquist frequency $f_N$ is:
\begin{equation}
    |H(f_N)| = \frac{1}{\sqrt{1 + (f_N/f_{3\,\mathrm{dB}})^2}}.
    \label{eq:nyquist_attenuation}
\end{equation}
For PAM-4 at a line rate of \SI{200}{\giga\bit\per\second}, the symbol rate is
$f_s = \SI{100}{\giga\baud}$ and the Nyquist frequency is $f_N = f_s/2 =
\SI{50}{\giga\hertz}$, giving $|H(\SI{50}{\giga\hertz})| = 1/\sqrt{1 +
(50/103)^2} \approx 0.90$, corresponding to less than \SI{1}{\decibel} of
Nyquist-frequency attenuation. For OOK at $f_s = \SI{100}{\giga\baud}$, the
Nyquist frequency coincides with $f_N = \SI{50}{\giga\hertz}$, yielding the
same attenuation. These values confirm that both formats are viable within the
passive $-3\,\mathrm{dB}$ bandwidth of the device without any electrical
equalisation, consistent with the approximation $f_N \leq 0.7\,f_{3\,\mathrm{dB}}$
cited in \cref{subsec:pd_format}. For a hypothetical upgrade to
\SI{400}{\giga\bit\per\second} PAM-4 (symbol rate $f_s = \SI{200}{\giga\baud}$,
$f_N = \SI{100}{\giga\hertz}$), the attenuation at the Nyquist frequency rises
to $|H(\SI{100}{\giga\hertz})| = 1/\sqrt{2} \approx -3\,\mathrm{dB}$—precisely
at the $-3\,\mathrm{dB}$ boundary—so that the signal power at the highest
temporal frequency component of the PAM-4 eye is halved relative to the DC
responsivity, requiring at minimum a single-tap linear equaliser in the receiver
DSP to compensate the Nyquist roll-off. The minimum receiver sensitivity at
\SI{400}{\giga\bit\per\second} PAM-4 is further degraded relative to
\cref{eq:sens_result} by the equalisation noise enhancement factor—the penalty
introduced by boosting the noise floor at high frequencies to flatten the
channel response—which for a first-order channel and a linear equaliser scales
as $\mathrm{NEF} \approx \sqrt{1 + (f_N/f_{3\,\mathrm{dB}})^2}$ and equals
$\sqrt{2} \approx \SI{1.5}{\decibel}$ at $f_N = f_{3\,\mathrm{dB}}$. This
noise enhancement penalty, which originates in the photodetector transfer
function and propagates through the equaliser into the receiver bit-error rate
budget, provides the direct quantitative connection between the device-level
$f_{3\,\mathrm{dB}}$ of \cref{eq:bw_total} and the system-level DSP complexity
discussed in the companion chapters~\cite{Yan2024}.


% =============================================================================
% SUGGESTED ADDITIONAL FIGURES (Only critical additions)
% =============================================================================

% [Suggested Figure – \cref{subsec:pd_compare} / Extension R6]
% Description : Three-panel comparison of the absorption fraction A(L) from
%   eq:abs_fraction (single-pass Beer-Lambert), the RCE absorption
%   A_RCE from eq:rce_absorption at the resonance condition (for R1=0.3,
%   R2=0.99), and the area-normalised dark current density J_dark as
%   functions of absorber thickness W_i (for the RCE case, L = W_i) from
%   10 nm to 800 nm. Overlay the transit-time-limited bandwidth f_tt(W_i)
%   on a secondary y-axis. All three absorption curves converge at large W_i;
%   at small W_i the RCE curve maintains A near unity while J_dark drops.
% Purpose : Visually establishes why RCE structures decouple the
%   absorption-thickness trade-off from the transit-time-bandwidth trade-off,
%   and quantifies the dark current floor reduction achievable by thinning
%   the absorber within a resonant architecture. Connects eq:rce_absorption to
%   eq:bw_tt and eq:jdark_srh in a single design-space diagram.

% [Suggested Figure – \cref{sec:pd_discussion} / Extension R9]
% Description : Contour plot of f_3dB (GHz) as a function of junction area A_j
%   (x-axis, 1 to 20 µm^2) and intrinsic layer thickness W_i (y-axis, 100 to
%   700 nm), computed from eq:bw_total with eq:bw_tt, eq:bw_rc, eq:cj, and
%   the U-shaped electrode R_s model from eq:rs_anode. Overlay the
%   Beer-Lambert absorption saturation contour A(L) = 0.90 (i.e. the minimum
%   L for 90% absorption at the device's A_j/W_m ratio), which bounds the
%   design space from below in A_j. Mark the present device operating point
%   (A_j = 6 µm^2, W_i = 350 nm) and the W_i* locus from eq:wi_optimal.
% Purpose : Provides a two-dimensional design map that consolidates all the
%   competing constraints (transit time, RC, absorption, dark current volume)
%   into a single diagram, making the optimality of the chosen geometry
%   visually evident and identifying the trajectory along which future scaling
%   should proceed.

% [Suggested Figure – \cref{subsec:pd_format} / Extension R10]
% Description : Normalised electrical power transfer function |H(f)|^2 (dB)
%   versus frequency from 1 to 250 GHz, with three curves: the present device
%   at f_3dB = 103 GHz; a hypothetical 50% R_s reduction device at ~130 GHz;
%   and the theoretical transit-time ceiling at f_tt = 51 GHz alone. Overlay
%   vertical dashed lines at the Nyquist frequencies for OOK/100G (50 GHz),
%   PAM4/200G (50 GHz), and PAM4/400G (100 GHz). Annotate the Nyquist
%   attenuation and the equalisation noise enhancement factor NEF at each
%   Nyquist point from eq:nyquist_attenuation.
% Purpose : Provides the quantitative link between the photodetector frequency
%   response and the DSP equalisation burden, translating device metrics
%   directly into system design requirements for the companion chapters.
% =============================================================================
